% Self-contained section body. The parent document supplies amsmath,
% amssymb, amsthm, and a proposition environment.
% Bibliography keys geco and tiny are already present in phase5 references.
% Primary sources checked: arXiv:1106.1622v1 and arXiv:2405.19816v2.
\section{Mathematical scope and the budget endpoint}
\label{sec:p6-theory}

The mathematical ingredients below are classical Gaussian moment
identities and elementary properties of the chosen evaluation endpoint.
They are not claimed as novel learning principles. Greedy low-rank
optimization and growth based on expressivity bottlenecks already have
established antecedents~\cite{geco,tiny}. Our question is empirical:
does the information supplied by a historical moment improve the
training-budget frontier once its computational costs are charged?

\subsection{Exact population risk of the trained predictor}

Let $X\sim N(0,I_d)$, $H(X)=XX^\top-I_d$, and
$q_A(X)=\langle A,H(X)\rangle$, where matrices are symmetric and
$\langle A,B\rangle=\operatorname{tr}(AB)$. The teacher and predictor are
\begin{equation}
 Y=q_{A_*}(X)+c_*+\xi,
 \qquad f_{B,c}(X)=q_B(X)+c,
 \qquad B=\sum_{j=1}^k a_jw_jw_j^\top,
 \label{eq:p6-model}
\end{equation}
where $\xi$ is independent of $X$, with mean zero and variance
$\sigma^2<\infty$. Our experiments set $c_*=0$ but fit $c$.

\begin{proposition}[Exact risk including the intercept]
\label{prop:p6-risk}
For every symmetric $B$ and scalar $c$,
\begin{equation}
 R(B,c):=\mathbb E[(Y-f_{B,c}(X))^2]
 =\sigma^2+2\|B-A_*\|_F^2+(c-c_*)^2.
 \label{eq:p6-risk}
\end{equation}
Consequently, the excess risk used in the experiments is
$\mathcal E(B,c)=2\|B-A_*\|_F^2+c^2$.
The same formula holds conditionally on a realized training history and
teacher, provided the test example is independent of that history.
\end{proposition}
\begin{proof}
The Gaussian second- and fourth-moment identities give
\[
 \mathbb E[X_iX_j]=\delta_{ij},\qquad
 \mathbb E[X_iX_jX_kX_l]
 =\delta_{ij}\delta_{kl}+\delta_{ik}\delta_{jl}
  +\delta_{il}\delta_{jk}.
\]
Thus $\mathbb E H_{ij}=0$ and
$\mathbb E[H_{ij}H_{kl}]=\delta_{ik}\delta_{jl}
+\delta_{il}\delta_{jk}$. Multiplying by $D_{ij}D_{kl}$ and summing
gives $\mathbb E[q_D(X)^2]=2\|D\|_F^2$ for symmetric $D$.
Take $D=A_*-B$. Since $\mathbb E q_D=0$, the cross term between
$q_D$ and $c_*-c$ vanishes. Noise independence and zero mean eliminate
its cross terms with both quantities. Expanding
$(q_D+c_*-c+\xi)^2$ proves~\eqref{eq:p6-risk}.
Conditioning on a saved predictor and its training history fixes $B,c$;
the independent test input retains the same Gaussian moments.
\end{proof}

The intercept term cannot be omitted: $B=A_*$ and $c=1$, for example,
have excess risk one when $c_*=0$. For each individual run, a separate
evaluator applies \eqref{eq:p6-risk} after its training trajectory has
been fixed. Development and confirmation have distinct roles. Before
observing any new development results, we froze a common search protocol:
all four methods are evaluated at learning rates
$\{0.0075,0.015,0.03\}$ and all five budgets on 16 independent
development teacher/data seeds. For each method, the mean exact
population AUC across these development seeds selects one learning
rate, which is then frozen before the 256-replicate confirmation
experiment. This is explicitly an oracle-based development selection;
its computational cost is reported separately. Confirmation teacher
matrices, population scores and oracle direction coefficients do not
enter the learner's proposal selection, optimization, stopping, or
hyperparameter selection. Exact evaluation removes test-input Monte
Carlo error in this synthetic model; it does not remove variation
across independently generated training problems.

\subsection{Function-preserving birth and the limits of direction quality}

Appending any finite direction $v$ with output coefficient
$a_{k+1}=0$ leaves both $B$ and $c$ unchanged, and hence preserves the
predictor for every input. Initializing the new Adam moment states to
zero also leaves the function unchanged. The subsequent optimization
trajectory is a separate matter: at birth the gradient with respect to
the new hidden vector is zero because it is multiplied by its output
coefficient, while the new output coefficient can have a nonzero
gradient.

For a unit direction $v$, set $E=A_*-B$ and hold all other parameters,
including the intercept, fixed. Proposition~\ref{prop:p6-risk} yields
\begin{equation}
 R(B+\alpha vv^\top,c)-R(B,c)
 =2\alpha^2-4\alpha v^\top Ev.
 \label{eq:p6-direction-gain}
\end{equation}
The oracle scalar coefficient is $\alpha=v^\top Ev$, with improvement
$2(v^\top Ev)^2$. Maximizing this quantity over unit directions gives
$2\|E\|_{\mathrm{op}}^2$. This is the improvement available from an
ideal scalar refit at a fixed predictor, not an improvement guaranteed
by the implemented factor Adam updates.

In particular, the earlier project's fresh-block quadratic analysis
concerns a prescribed update of the form
\[
 B^+=B+\frac{1}{2n}\sum_{i=1}^n
       \{Y_i-q_B(X_i)\}H(X_i),
\]
under its stated zero-intercept model and conditional independence
assumptions. Factor Adam instead updates $(w_j,a_j,c)$ using empirical
gradients, accumulated coordinatewise moments, and clipping. It may
reuse retained examples and historical summaries. Neither its update
map nor its sample dependence matches that matrix recurrence, so the
contraction theorem does not supply a convergence rate or cost bound
for the present experiment. Even the exact empirical residual matrix,
when the predictor has an intercept, is
\begin{align*}
 G_n(B,c)&=S_n-T_n(B)-\tfrac12c\overline H_n,
 &S_n&=\frac1{2n}\sum_iY_iH_i,\\
 T_n(B)&=\frac1{2n}\sum_iq_B(X_i)H_i,
 &\overline H_n&=\frac1n\sum_iH_i.
\end{align*}
The historical proxy $S_n-B$ uses the known population Gaussian
operator; it is not this exact empirical residual.

\subsection{A fixed log-budget area from five completed runs}

The five preregistered per-arrival budget caps, in the declared work
units, are
\[
 (b_1,b_2,b_3,b_4,b_5)
 =10^6(3.2,5,6.4,8,10).
\]
Every run processes seven arrivals under the same data-availability
rule. Let $e_{mrj}$ be the final excess risk of method $m$, replicate
$r$, and budget setting $b_j$, after all seven arrivals. Each entry is
obtained from its own completed training run. These five entries are
not checkpoints extracted from one run with the largest budget.

Write $x_j=\log b_j$ and $L=x_5-x_1>0$. Define the normalized
log-budget trapezoid area
\begin{equation}
 \operatorname{AUC}_{mr}
 =\frac1L\sum_{j=1}^{4}(x_{j+1}-x_j)
                   \frac{e_{mrj}+e_{mr,j+1}}2
 =\sum_{j=1}^5\omega_j e_{mrj},
 \label{eq:p6-auc}
\end{equation}
where
\[
 \omega_1=\frac{x_2-x_1}{2L},\qquad
 \omega_5=\frac{x_5-x_4}{2L},\qquad
 \omega_j=\frac{x_{j+1}-x_{j-1}}{2L}\quad(2\le j\le4).
\]
The weights are positive, sum to one, and are fixed independently of
all outcomes. Only budget is logarithmically transformed; the ordinate
is excess risk itself. Lower area is better. Replacing $b_j$ by its
seven-arrival cap $7b_j$ leaves all weights unchanged, as does changing
the logarithm base.

Equation~\eqref{eq:p6-auc} is an integral of the piecewise-linear
interpolant in log budget, not an assumption that the unobserved
performance curve is linear. The budget axis is the assigned cap;
actual charged work, unused terminal budget and wall time are reported
separately. The area therefore summarizes the specified allocation
protocol and cost ledger, and does not by itself establish hardware
runtime dominance. The primary sample comprises 256 independent
teacher/data replicates, with pairing across methods and budget
settings within each replicate. The five budgets do not multiply the
number of independent inference units by five.

\subsection{Meaning of the three-comparison intersection--union test}

Let $h$ denote historical-moment growth, and let $m_1,m_2,m_3$ denote
the three preregistered primary comparators. For $\ell=1,2,3$, define
the paired area difference and its mean over the specified distribution
of independently generated problems by
\[
 D_{r\ell}=\operatorname{AUC}_{hr}
                -\operatorname{AUC}_{m_\ell r},
 \qquad \mu_\ell=\mathbb E[D_{r\ell}].
\]
The three-comparison intersection--union test (IUT3) addresses
\begin{equation}
 H_0:\ \bigcup_{\ell=1}^3\{\mu_\ell\ge0\},
 \qquad
 H_1:\ \bigcap_{\ell=1}^3\{\mu_\ell<0\}.
 \label{eq:p6-iut}
\end{equation}
If $p_\ell$ is a valid one-sided $p$-value for the corresponding
component null, reject~\eqref{eq:p6-iut} at level $\alpha$ only when
all three $p_\ell\le\alpha$, equivalently when
$p_{\rm IUT}=\max_\ell p_\ell\le\alpha$.
Indeed, under the union null at least one component null, say
$\ell_0$, is true, and
\[
 \mathbb P(p_{\rm IUT}\le\alpha)
 \le\mathbb P(p_{\ell_0}\le\alpha)\le\alpha.
\]
This argument requires no independence between the three comparisons
and no Bonferroni division of $\alpha$ for this single conjunction
claim. Its calibration still depends on the validity of each chosen
component test; a finite-sample exact claim cannot be inferred merely
from having 256 replicates. Additional families of claims or
outcome-dependent selection of comparators require their own handling.

Rejection supports lower mean area against every one of the three
specified comparators under this protocol. It does not imply lower
risk at every budget, success on every seed, superiority to untested
methods, or optimality among architectures. Failure to reject does
not establish equivalence, absence of a directional effect, or absence
of an advantage at some budgets. A superiority test has no prespecified
equivalence margin. Together, the risk identity, direction diagnostic
and IUT3 interpretation provide no uniform optimal cost bound for
factor Adam or autonomous network growth.
